% =================================================================
% ARCHIVO DE DEFINICIÓN DE GLOSARIO DE TÉRMINOS (glosario.tex)
% =================================================================

% Para añadir un término al glosario, usa el siguiente comando:
% \newglossaryentry{etiqueta_corta}{propiedades}

% ------------------------------------------------------------------------------------------
% PROPIEDADES MÍNIMAS REQUERIDAS:
% - {etiqueta_corta}: El ID interno para hacer referencia (ej. 'robot', 'latex').
% - name={Nombre}: El nombre que aparecerá en el glosario (ej. 'Robot', 'Latex').
%                  ¡Debe ir con la primera letra en mayúscula para mantener el formato UNEFA!
% - description={...}: La definición completa del término.
% ------------------------------------------------------------------------------------------

% USO EN EL DOCUMENTO:
% - Los términos SÓLO aparecerán en el glosario si son llamados al menos una vez en el cuerpo.
% - \gls{etiqueta_corta}: Inserta el término en minúsculas (ej. robot).
% - \Gls{etiqueta_corta}: Inserta el término con la primera letra en MAYÚSCULA (ej. Robot).
% - \glspl{etiqueta_corta} / \Glspl{etiqueta_corta}: Versiones PLURALES.
% ------------------------------------------------------------------------------------------

% EJEMPLO:
% \newglossaryentry{arduino}
% {
%     name={Arduino},
%     description={Arduino es una tecnología de hardware de código abierto...}
% }
% 
% Uso en el texto: "La placa de desarrollo es un \gls{arduino}"

% Para más información, consulte info_glosario.md


\newglossaryentry{latex}
{
    name=Latex,
    description={Es un lenguaje de enmaquetado especialmente útil
    para documentos científicos}
}

\newglossaryentry{maths}
{
    name=Matemáticas,
    description={es lo que hacen los matemáticos}
}

\newglossaryentry{vacio}
{
    name=Vacío,
    description={Término que nunca se empleó, y por lo tanto no debería
    aparecer en el glosario}
}
